\documentclass[11pt,a4paper]{article}
\usepackage[margin=22mm]{geometry}
\usepackage{amsmath,amssymb,booktabs,microtype}
\usepackage[hidelinks]{hyperref}
\newcommand{\dd}{\mathrm d}
\newcommand{\one}{\mathbf 1}
\title{Route G-R6: Separating Thermal Response\\from a Gravitational Clock Baseline}
\author{A fixed-input two-temperature, two-potential control}
\date{29 September 2026}
\begin{document}
\maketitle
\begin{abstract}
Retain G-R4's E2/E3 thermal action and inputs and G-R5's external E3
reference. Add an independently calibrated gravitational-potential
contrast. A four-setting logarithmic comparison separates thermal,
gravitational and interaction responses without fitting a time law.
Matched local radiation states cancel even uncertain fixed thermal
shifts in the height contrast. Uncontrolled height-correlated biases
remain degenerate with anomalous redshift. This is a theory baseline,
not a new experiment, an empirical limit, or a derivation of spacetime.
\end{abstract}

\section{Geometry acts on clock rate; radiation acts on local spectrum}
Use the test-clock approximation: tidal effects across the ion and
heater-induced changes in the background are negligible or independently
bounded. Assume a static background with
\begin{equation}
 \dd s^2=-N(\mathbf x)^2c^2\dd t^2+h_{ab}\dd x^a\dd x^b .
\end{equation}
For a clock held at rest, $\dd\tau=N\dd t$. A local transition with
proper frequency $\nu_i^{\rm loc}$ accumulates phase
$\dd\varphi_i=2\pi\nu_i^{\rm loc}\dd\tau$. Stationary phase transfer
to a reference at $R$ therefore gives
\begin{equation}\label{eq:ratio}
 \frac{\nu_{i,\mathrm{received}}^S}{\nu_3^R}
 =\frac{\nu_i^0}{\nu_3^0}\frac{N_S}{N_R}
   \frac{1+b_i^S}{1+b_3^R}\exp(A_i^S-A_3^R+\ell_i).
\end{equation}
Here $A$ denotes other local log-frequency shifts and $\ell_i$
nonideal transfer corrections. They are not silently zero in data.
G-R4's local dipole response is unchanged:
\begin{equation}
 b_i^X=-\frac{1}{2h\epsilon_0\nu_i^0}\mathcal P
 \int_0^\infty\Delta\alpha_i(\omega)u_\omega^X\dd\omega,
 \qquad B_i^X=\ln(1+b_i^X).
\end{equation}
For the retained ideal static approximation,
$b_i(T)=\beta_i a_{\rm rad}T^4$ and
$\beta_i=-\Delta\alpha_i/(2h\epsilon_0\nu_i^0)$.
All temperatures and spectra are local proper quantities.

Normalize Eq.~\eqref{eq:ratio} by $\nu_i^0/\nu_3^0$ and take its log:
\begin{equation}\label{eq:q}
 q_i=G+B_i^S-B_3^R+A_i^S-A_3^R+\ell_i+n_i,\quad
 G=\ln(N_S/N_R).
\end{equation}
The universal gravitational factor is now specified by geometry,
not a fitted replacement for G-R5's unidentified common factor.
To first weak-field order,
$\Delta G=\Delta\Phi/c^2$ for stationary nonrotating coordinates.
For moving clocks an additional $-\Delta(v^2)/(2c^2)$ appears.
This is a restricted redshift/local-position-invariance baseline,
not a test of every classical or quantum EEP component~\cite{zych}.

\newpage
\section{Four settings, three contrasts}
Use the same target ion and enclosure, with E2/E3 interleaved at each
settled plateau. Keep the independent E3 reference at its fixed site.
Let $a=0,1$ label local temperature and $b=0,1$ label potential.
The column order is $(00,10,01,11)$. After evaluating reference,
transfer and non-BBR shifts, the ideal model is
\begin{equation}\label{eq:separable}
 q_{i,ab}=q_{i,00}+a\theta_i+b\gamma,\qquad
 \theta_i=B_i(T_1)-B_i(T_0),\quad
 \gamma=\ln\frac{N_1}{N_0}.
\end{equation}
This holds for any fixed local $B_i(T)$, not only its static $T^4$
approximation, provided the same local spectrum and apparatus
conditions are reproduced at both potential settings.
Define
\begin{align}
 \mathcal T_i&=\tfrac12(q_{i,10}-q_{i,00}+q_{i,11}-q_{i,01})
                =\theta_i,\\
 \mathcal G_i&=\tfrac12(q_{i,01}-q_{i,00}+q_{i,11}-q_{i,10})
                =\gamma,\\
 \mathcal I_i&=q_{i,11}-q_{i,10}-q_{i,01}+q_{i,00}=0.
\end{align}
Thus the null predictions are
\begin{equation}\label{eq:null}
 \mathcal G_2=\mathcal G_3=\gamma,\qquad
 \mathcal G_3-\mathcal G_2=0,\qquad \mathcal I_2=\mathcal I_3=0.
\end{equation}
The internal E3/E2 ratio cancels gravity: it cannot by itself measure
the common redshift. The outside reference is essential.

\paragraph{One illustrative calculation, not observations.}
Retain the pinned G-R4 polarizabilities, $T_0=300$ K and $T_1=310$ K.
Take $\Delta h=1$ m and conventional $g_*=9.80665$ m\,s$^{-2}$ merely
as a scale illustration, not measured laboratory geodesy. Then
$\gamma\simeq g_*\Delta h/c^2=1.0911\,10^{-16}$:
\begin{center}
\begin{tabular}{lrr}
\toprule
Setting & $(q_2-q_{2,00})/10^{-16}$ & $(q_3-q_{3,00})/10^{-16}$\\
\midrule
$T_0,\Phi_0$ & $0$ & $0$\\
$T_1,\Phi_0$ & $-0.7337$ & $-0.1012$\\
$T_0,\Phi_1$ & $1.0911$ & $1.0911$\\
$T_1,\Phi_1$ & $0.3574$ & $0.9899$\\
\bottomrule
\end{tabular}
\end{center}
The thermal E3/E2 contrast is $6.325\,10^{-17}$; the gravitational
E3/E2 contrast is zero. All entries are model predictions.

\paragraph{Why use logarithms?}
For the baseline-normalized raw ratio $r_{ab}=\exp(a\theta_i+b\gamma)$,
the raw cross contrast is
\begin{equation}
 r_{11}-r_{10}-r_{01}+r_{00}
 =(e^{\theta_i}-1)(e^\gamma-1),
\end{equation}
not zero. This ordinary multiplicative product is not new physics.
Logarithms make the separability null exact within the stated model.

\newpage
\section{What cancels, what does not, and identifiability}
An arbitrary but height-independent $B_i(T_a)$ cancels in
$\mathcal G_i$. Consequently the old polarizability uncertainty need
not impose its full G-R5 thermal-composite floor on the matched-bath
height test. This is not permission to neglect radiation matching.
If the higher site instead has $T_a+\delta T_a$, then
\begin{equation}\label{eq:leak}
 \delta\mathcal G_i^{\rm BBR}
 \simeq\frac12\sum_{a=0}^1 B_i'(T_a)\delta T_a,\qquad
 B_i'(T)=\frac{4b_i(T)}{T[1+b_i(T)]}
\end{equation}
in the retained static model. The exact formula is one half the sum
of $B_i(T_a+\delta T_a)-B_i(T_a)$.
An illustrative $+0.1$ K mismatch at both upper-site setpoints leaks
$-7.345\,10^{-19}$ for E2 and $-1.013\,10^{-19}$ for E3.
Dynamic, anisotropic and geometry changes require their own responses;
matching a wall thermometer alone does not establish matching spectra.

G-R5's composite weights
\begin{equation}
 w_2=\frac{\beta_3}{\beta_3-\beta_2},\quad
 w_3=-\frac{\beta_2}{\beta_3-\beta_2},\quad
 (w_2,w_3)\simeq(-0.16004,1.16004)
\end{equation}
remain a leading thermal-rejection diagnostic. They obey
$w_2+w_3=1$ and $w_2\beta_2+w_3\beta_3=0$, so the common height
contrast is retained. Their use is not necessary for the exact
matched-bath cancellation above.

For interpretation only, parametrize a possible anomalous height
response by $\alpha_i$:
\begin{equation}
 \mathcal G_i=(1+\alpha_i)\gamma+c_h+\varepsilon_i .
\end{equation}
The baseline is $\alpha_i=0$, not a fit. A shared height-correlated
reference/link bias $c_h$ then gives
\begin{align}
 D_h&=\mathcal G_3-\mathcal G_2
       =(\alpha_3-\alpha_2)\gamma+\varepsilon_3-\varepsilon_2,\\
 C_h-\gamma&=(w_2\alpha_2+w_3\alpha_3)\gamma
              +c_h+w_2\varepsilon_2+w_3\varepsilon_3.
\end{align}
For $\gamma\ne0$, shifting both $\alpha_i$ by $\lambda$ and $c_h$
by $-\lambda\gamma$ leaves the observations unchanged. A common
anomaly is therefore unidentifiable without independent bias bounds.
Differential biases can also mimic $D_h$; a nonzero interaction
contrast only falsifies the controlled separability model, not EEP
by itself.

Order the eight channel means as E2's four cells followed by E3's.
For any contrast vector $a$, its variance is $a^TVa$. A shared,
cell-varying reference contribution has covariance
$V_R=\one_2\one_2^T\otimes K_R$. It cancels for
$a_D=(-1,1)\otimes(-1,-1,1,1)/2$, but not generally for
$a_C=w\otimes(-1,-1,1,1)/2$. A constant reference offset cancels
from both. Matched effective sampling windows are required.

\newpage
\section{Independent controls and a bounded empirical handoff}
\paragraph{Potential.}
Measure $\Delta\Phi$ independently at the ion positions using
leveling/gravimetry/geodesy with covariance. A portable-clock campaign
with independent geodetic comparison~\cite{grotti} reports the magnitude
$3915.88(0.30)$ m$^2$\,s$^{-2}$. This gives a source-backed baseline
$|\Delta G|\simeq4.3570\,10^{-14}$ with propagated potential uncertainty
$3.34\,10^{-18}$. We use the geodetic value, not the clock-inferred
potential. The published clocks were Sr: these are not our E2/E3
four-setting data or a calibration of our apparatus.
For this upward-positive convention,
$\Delta\Phi=\int_{h_0}^{h_1}g(h)\dd h$ locally.
If a terrestrial geopotential already includes centrifugal effects,
do not add the corresponding rotational kinetic term again.
Specify sign and transfer convention; stationary-plateau frequency
comparisons are not transported-clock accumulated-phase measurements.
A $10^{-18}$ absolute redshift error corresponds to
$\sigma_{\Delta\Phi}\simeq0.0899$ m$^2$\,s$^{-2}$, or about
$0.92$ cm at $g_*$, for this contribution alone. No such accuracy
is claimed. In a residual $r=\mathcal G-\gamma_{\rm cal}$, propagate
$\operatorname{Var}r=\operatorname{Var}\mathcal G+
\operatorname{Var}\gamma_{\rm cal}-2\operatorname{Cov}(\mathcal G,\gamma_{\rm cal})$.

\paragraph{Temperature and reference.}
Use independent local radiometry, with sensor calibration, gradients,
emissivity/view factors and reference monitoring at every plateau.
Do not infer the bath from the tested frequency law, or infer potential
from the same clock ratios and then call the prediction verified.
Radiometry/correction covariance enters $V$ through its response
Jacobian. Actual measured settings, not ideal labels, enter the model.

\paragraph{Equilibrium is a different boundary condition.}
Independently thermostatted local setpoints need not form one global
passive equilibrium bath. In a static globally equilibrated radiation
field, conservation of Killing energy gives $E_\infty=NE_{\rm loc}$;
the Bose factor written locally therefore has
\begin{equation}
 e^{E_\infty/(k_BT_\infty)}
 =e^{E_{\rm loc}/(k_BT_{\rm loc})},\qquad
 T_{\rm loc}N=T_\infty .
\end{equation}
This is the Tolman--Ehrenfest relation~\cite{tolman,thermal}.
Do not impose it simultaneously with freely chosen local temperatures.
At 300 K across the illustrative 1 m separation it predicts only
$\Delta T\simeq-300\gamma=-3.27\,10^{-14}$ K.
This observation neither changes the imposed 10 K control nor requires
resolving that tiny equilibrium gradient.

\paragraph{Acquisition and stopping rule.}
One optional symmetric order is $(00,10,11,01,01,11,10,00)$ at
centers $(-7,-5,-3,-1,1,3,5,7)\tau$. Pair averaging cancels ideal
linear drift in all three contrasts, but quadratic drift survives:
its $(\mathcal T,\mathcal G,\mathcal I)$ moments are
$(-8,-32,32)\tau^2$. State-locked bias also survives.
Actual settling/windows, return comparisons, reversed orders,
transport-systematic monitoring and repeated blocks remain necessary.
Four means do not determine statistical significance.
The input card records the missing acquisition fields. No paired data,
new experiment or lab contact is supplied; no time law or Route-F gate
is promoted. The baseline is complete within its assumptions.

\begin{thebibliography}{9}\small
\bibitem{zych} M. Zych and C. Brukner,
\href{https://doi.org/10.1038/s41567-018-0197-6}{Nature Physics 14, 1027 (2018)}.
\bibitem{grotti} J. Grotti et al.,
\href{https://doi.org/10.1103/PhysRevApplied.21.L061001}{Physical Review Applied 21, L061001 (2024)}.
\bibitem{tolman} R. C. Tolman and P. Ehrenfest,
\href{https://doi.org/10.1103/PhysRev.36.1791}{Physical Review 36, 1791 (1930)}.
\bibitem{thermal} N. Karolinski and V. Faraoni,
\href{https://doi.org/10.1140/epjc/s10052-024-13632-6}{European Physical Journal C 84, 1248 (2024)}, section 1.
\end{thebibliography}
\end{document}
